% TOP_PROBLEM: {"id":30007639,"problem_number":"LOCAL-30007639","title":"Platform-worker protection and flexibility","category_id":21,"category":{"id":21,"name":"economics","display_name":"Economics","description":"Open economic theory statements, published research agendas, and proposed scoped specifications; question provenance and historical priority are qualified per record.","slug":"economics","order_index":21,"created_at":"2026-10-08T00:00:00Z"},"status":"provenance_pending","external_url":null,"legacy_ids":[],"published":false,"statement":"In the explicitly specified setting: Which combinations of employment status, portable benefits, pay floors, and scheduling flexibility improve platform-worker welfare without undermining access to work? The mathematical statement above fixes the target, assumptions and scope of this catalogue formulation.","economics":{"original_id":128,"field":"Labor markets and work","kind":"design","formalization_basis":"proposed_specification","source_status":"not_verified","priority_status":"not_established","priority_note":"No qualifying problem statement was directly verified, so historical priority cannot be assigned.","earliest_verified_reference":null,"current_reference":{"authors":["Christopher T. Stanton","Catherine Thomas"],"title":"Who Benefits from Online Gig Economy Platforms?","year":2025,"url":"https://www.hbs.edu/ris/Publication%20Files/StantonThomas_Submission3_c26466ef-fcf2-44c6-a356-c92d042113d2.pdf","doi":"10.1257/aer.20221189","locator":"November 2024 manuscript, section 5, printed pp. 29–31; section 6, printed pp. 31–32","evidence_summary":"The paper evaluates particular taxes and wage floors; no explicit open question about jointly optimal portable benefits and status was verified."},"formalization_note":"This is a proposed scoped research specification. No primary passage explicitly stating this entire remaining question as open was verified. The supporting literature may answer important subquestions; this entry must not be advertised as a verified formal open problem."},"statusQualification":"Provenance pending: an explicit published statement of this exact open question was not verified. This is a catalogue-authored candidate specification, not a certified open conjecture. This is a proposed scoped research specification. No primary passage explicitly stating this entire remaining question as open was verified. The supporting literature may answer important subquestions; this entry must not be advertised as a verified formal open problem.","statusReviewedAt":"2026-10-08"}
% FIRST_POSTED: 2026-10-08
% ENTRY_KIND: statement_only
% !TeX program = lualatex
\documentclass[11pt]{article}
\usepackage{fontspec}
\usepackage[a4paper,margin=25mm]{geometry}
\usepackage{amsmath,amssymb,mathtools}
\usepackage{xurl}
\usepackage[hidelinks]{hyperref}
\newcommand{\sref}[2]{\hyperlink{source:#1:#2}{[S#2]}}
\setlength{\emergencystretch}{3em}
\begin{document}
% problemId:30007639
\section*{Platform-worker protection and flexibility}\label{30007639}
\subsection{Definitions and mathematical statement}
Fix a platform labor market and a policy menu \(\mathcal P\). Policy \(p\) specifies employment classification, an earnings floor, benefit contributions that follow workers across employers, and enforcement. Workers choose participation and hours; buyers choose postings and hires; platforms may change fees or allocation rules. Let \(U_i(p)\) be worker \(i\)'s expected discounted utility, including income risk, benefits, work flexibility, and application costs. Let \(N(p)\) be the number of workers obtaining paid tasks and \(C(p)\) total implementation cost. For explicit welfare weights \(\omega_i\), define \(S(p)=E[\sum_i\omega_iU_i(p)]\). Determine the Pareto frontier of \((S(p),N(p),-C(p))\), including unsuccessful applicants rather than conditioning solely on workers hired. Compare the frontier across plausible demand elasticities and outside employment options. Empirical identification requires policy or contract variation and a model of buyer, worker, and platform responses. A portable-benefits contribution must be traced through fees, wages, prices, and hiring. This is a proposed design problem motivated by existing platform evidence; the cited study does not establish that the entire regulatory menu remains an explicitly stated open problem, and its particular policy results must be preserved.

\par\textbf{Formulation and scope.} This is a proposed scoped research specification. No primary passage explicitly stating this entire remaining question as open was verified. The supporting literature may answer important subquestions; this entry must not be advertised as a verified formal open problem.

\par\textbf{Open-question provenance.} An explicit published open-question statement was not verified. Sources below are supporting literature and must not be treated as a first listing.

\par\textbf{Historical priority.} No qualifying problem statement was directly verified, so historical priority cannot be assigned.

\subsection{Short English statement}
In the explicitly specified setting: Which combinations of employment status, portable benefits, pay floors, and scheduling flexibility improve platform-worker welfare without undermining access to work? The mathematical statement above fixes the target, assumptions and scope of this catalogue formulation.

\subsection{Sources}
\begin{itemize}\raggedright
\item[\textnormal{[S1]}]\hypertarget{source:30007639:1}{} Christopher T. Stanton, Catherine Thomas. 2025. Who Benefits from Online Gig Economy Platforms?. November 2024 manuscript, section 5, printed pp. 29–31; section 6, printed pp. 31–32.
\url{https://www.hbs.edu/ris/Publication%20Files/StantonThomas_Submission3_c26466ef-fcf2-44c6-a356-c92d042113d2.pdf}.
DOI: 10.1257/aer.20221189.
{\small\textit{Supporting or current literature; see evidence qualification. The paper evaluates particular taxes and wage floors; no explicit open question about jointly optimal portable benefits and status was verified.}}
\end{itemize}
\end{document}
